\usepackage[utf8]{inputenc}
\usepackage[T1]{fontenc}
\usepackage[french]{babel}
\usepackage[left=3cm, right=2cm, top=2cm, bottom=2cm]{geometry} 
\usepackage{setspace}
\onehalfspacing 
\usepackage{fancyhdr}
\pagestyle{fancy}
\fancyhf{} 
\fancyhead[L]{\leftmark} 
\fancyfoot[C]{\thepage} 
\renewcommand{\headrulewidth}{0.4pt}
\renewcommand{\footrulewidth}{0.4pt}
\usepackage{graphicx} 
\usepackage{caption}
\usepackage{subcaption}
\usepackage{booktabs} 
\usepackage{array}
\usepackage{float}
\usepackage{tikz} 
\usepackage{pgfplots}
\pgfplotsset{compat=1.18}
\usepackage{listings}
\usepackage{xcolor}

\definecolor{codegreen}{rgb}{0,0.6,0}
\definecolor{codegray}{rgb}{0.5,0.5,0.5}
\definecolor{codepurple}{rgb}{0.58,0,0.82}
\definecolor{backcolour}{rgb}{0.96,0.96,0.96}

\lstdefinestyle{mystyle}{
    backgroundcolor=\color{backcolour},   
    commentstyle=\color{codegreen},
    keywordstyle=\color{blue}\bfseries,
    numberstyle=\tiny\color{codegray},
    stringstyle=\color{codepurple},
    basicstyle=\ttfamily\footnotesize,
    breakatwhitespace=false,         
    breaklines=true,                 
    captionpos=b,                    
    keepspaces=true,                 
    numbers=left,                    
    numbersep=5pt,                  
    showspaces=false,                
    showstringspaces=false,
    showtabs=false,                  
    tabsize=4,
    frame=single,
    rulecolor=\color{lightgray}
}
\lstset{style=mystyle}

\usepackage[backend=biber, style=ieee]{biblatex}
\addbibresource{references.bib} 
\usepackage{acronym}
\usepackage{lettrine} % Pour marquer le début d'un chapitre (lettrine)
\usepackage[hidelinks, colorlinks=true, linkcolor=black, citecolor=blue, urlcolor=blue]{hyperref}
%====================================================
% Packages recommandés pour tableaux modernes
%====================================================
\usepackage{array}
\usepackage{booktabs}
\usepackage{longtable}
\usepackage{tabularx}
\usepackage[table]{xcolor}
\usepackage{multirow}
\usepackage{makecell}
\usepackage{pifont}
\usepackage{enumitem}
\usepackage{url}
\usepackage{hyperref}

%====================================================
% Commandes visuelles
%====================================================
\newcommand{\cmark}{\ding{51}}
\newcommand{\xmark}{\ding{55}}
\newcommand{\pmark}{\ding{115}}

\definecolor{HeaderBlue}{HTML}{1E3A8A}
\definecolor{LightBlue}{HTML}{EFF6FF}
\definecolor{LightGreen}{HTML}{ECFDF5}
\definecolor{LightOrange}{HTML}{FFF7ED}
\definecolor{LightPurple}{HTML}{F5F3FF}
\definecolor{LightGray}{HTML}{F8FAFC}
\definecolor{SoftRed}{HTML}{FEF2F2}

\newcolumntype{Y}{>{\raggedright\arraybackslash}X}
\newcolumntype{C}{>{\centering\arraybackslash}p{1.8cm}}

%====================================================
% Tableaux modernes et compacts
%====================================================
\usepackage{array}
\usepackage{booktabs}
\usepackage{tabularx}
\usepackage{longtable}
\usepackage[table]{xcolor}
\usepackage{makecell}
\usepackage{ragged2e}
\usepackage{pifont}

\definecolor{HeaderBlue}{HTML}{1E3A8A}
\definecolor{LightGreen}{HTML}{ECFDF5}
\definecolor{LightOrange}{HTML}{FFF7ED}
\definecolor{LightPurple}{HTML}{F5F3FF}
\definecolor{LightBlue}{HTML}{EFF6FF}
\definecolor{SoftRed}{HTML}{FEF2F2}
\definecolor{LightGray}{HTML}{F8FAFC}

\newcolumntype{Y}{>{\RaggedRight\arraybackslash}X}
\newcolumntype{C}{>{\centering\arraybackslash}p{1.35cm}}

\newcommand{\oui}{\cellcolor{LightGreen}Oui}
\newcommand{\non}{\cellcolor{SoftRed}Non}
\newcommand{\partiel}{\cellcolor{LightOrange}Partiel}
\newcommand{\prep}{\cellcolor{LightPurple}Préparé}
\newcommand{\faible}{\cellcolor{SoftRed}Faible}
\newcommand{\moyen}{\cellcolor{LightOrange}Moyen}
\newcommand{\eleve}{\cellcolor{LightGreen}Élevé}